\hmsubsection{Performance Evaluation}{}{}

This section evaluates the performance of the key subsystems
developed for the sorting robot, covering the inverse kinematics
solver, the sag compensation model, and the adaptive grip
mechanism. Each subsystem is assessed against its design
requirements using measurements collected during development and
the automated test suite.

\hmsubsubsection{Inverse Kinematics Solver}{UMA}{}

The geometric inverse kinematics solver produces a solution in
less than 1~ms on the Raspberry Pi~5 host, which is well below
the time required for the physical motors to reach their goal
positions (typically 1--2~s). The solver's computational cost is
therefore negligible in the overall cycle time, and the dominant
bottleneck is the mechanical motion of the arm.

The solver's positional accuracy was validated through automated
forward-kinematics round-trip tests. For every target position
in a swept grid spanning the reachable workspace, the solver
computes motor positions, which are then converted back to
Cartesian coordinates using the forward kinematics function.
The maximum round-trip error across all tested positions was
below 0.1~mm, confirming that the geometric derivation and the
step-to-angle conversion are internally consistent.

A symmetry test further verified the solver by comparing
solutions for mirror targets $(x, +y, z)$ and $(x, -y, z)$.
In all cases, the shoulder, elbow, and wrist angles were
identical, and the base angles summed to exactly 4096~steps,
confirming the absence of sign errors in the derivation.

\hmsubsubsection{Sag Compensation Accuracy}{UMA}{}

The polynomial sag model achieved a root mean square error of
0.114~cm across the six calibration points. However, as
described in Section~8.7.3, the sag model is not exercised in
the deployed configuration because the touch calibration
procedure provides per-point surface heights that are more
accurate than the reach-based regression. The touch calibration
heights embed the arm's real-world droop empirically at each
calibrated distance, which eliminates the sag model's
assumption that droop depends only on horizontal reach and not
on base angle or arm configuration.

\hmsubsubsection{Adaptive Grip Reliability}{UMA}{}

The adaptive grip algorithm uses the Dynamixel motor's present
load and present current registers to detect whether a ball is
held. During testing with 50~mm balls at five workspace
positions, the grip verification correctly rejected empty grips
(claw closed on air) in every observed case, and false
negatives (ball present but grip rejected) were rare. The
two-tier verification threshold, with a 30-step position gate
for strong contact and a 5-step gate with sensor confirmation
for marginal contact, proved effective at distinguishing
between a gripped ball and an empty close. The 50~mA current
cap was low enough to detect resistance reliably while
protecting the 3D-printed claw from excessive force.
